\documentclass[UTF8,zihao=-4,oneside,fontset=fandol]{ctexrep}

\usepackage[a4paper,margin=2.35cm]{geometry}
\usepackage{amsmath,amssymb,amsthm,mathtools,bm}
\usepackage{booktabs,tabularx,array,multirow}
\usepackage{enumitem}
\usepackage{xcolor}
\usepackage{hyperref}
\usepackage{fancyhdr}
\usepackage{setspace}
\usepackage{listings}
\usepackage{tikz}
\usepackage[most]{tcolorbox}
\usetikzlibrary{arrows.meta,positioning,trees,calc}

\definecolor{DeepBlue}{HTML}{1F4E79}
\definecolor{DeepGreen}{HTML}{2E6B4F}
\definecolor{WarmRed}{HTML}{A33A32}
\definecolor{SoftGray}{HTML}{F4F5F6}

\hypersetup{
  colorlinks=true,
  linkcolor=DeepBlue,
  urlcolor=DeepBlue,
  citecolor=DeepGreen,
  pdftitle={信息论自学讲义},
  pdfauthor={RUCmonk2}
}

\setstretch{1.22}
\setlength{\parindent}{2em}
\setlength{\parskip}{0.25em}
\setlist{itemsep=0.25em,topsep=0.35em}

\pagestyle{fancy}
\fancyhf{}
\fancyhead[L]{信息论自学讲义}
\fancyhead[R]{\leftmark}
\fancyfoot[C]{\thepage}
\setlength{\headheight}{14pt}

\newtheorem{definition}{定义}[chapter]
\newtheorem{theorem}[definition]{定理}
\newtheorem{proposition}[definition]{命题}
\newtheorem{lemma}[definition]{引理}
\newtheorem{corollary}[definition]{推论}
\theoremstyle{remark}
\newtheorem{example}[definition]{例}
\newtheorem{remark}[definition]{说明}

\newtcolorbox{intuitionbox}{
  colback=DeepBlue!4,colframe=DeepBlue!70!black,
  title=直觉,breakable,boxrule=0.7pt,arc=1mm
}
\newtcolorbox{warningbox}{
  colback=WarmRed!4,colframe=WarmRed!75!black,
  title=容易混淆,breakable,boxrule=0.7pt,arc=1mm
}
\newtcolorbox{aibox}{
  colback=DeepGreen!4,colframe=DeepGreen!75!black,
  title=与 AI4Math / AI4TCS 的连接,breakable,boxrule=0.7pt,arc=1mm
}
\newtcolorbox{checkpointbox}{
  colback=SoftGray,colframe=black!50,
  title=本节验收,breakable,boxrule=0.6pt,arc=1mm
}

\lstset{
  basicstyle=\ttfamily\small,
  breaklines=true,
  columns=fullflexible,
  frame=single,
  rulecolor=\color{black!25},
  backgroundcolor=\color{black!2},
  xleftmargin=0.6em,
  xrightmargin=0.6em,
  showstringspaces=false
}

\newcommand{\E}{\mathbb{E}}
\newcommand{\Var}{\operatorname{Var}}
\newcommand{\Cov}{\operatorname{Cov}}
\newcommand{\Prb}{\mathbb{P}}
\newcommand{\KL}{D_{\mathrm{KL}}}
\newcommand{\ind}{\mathbf{1}}
\newcommand{\bits}{\text{ bits}}
\DeclareMathOperator*{\argmax}{arg\,max}

\title{\heiti 信息论自学讲义\\[0.5em]
\Large 从熵、编码到信道容量与可验证算法}
\author{RUCmonk2}
\date{2026 年 8 月}

\begin{document}

\hypersetup{pageanchor=false}
\maketitle
\hypersetup{pageanchor=true}

\chapter*{如何使用这本讲义}
\addcontentsline{toc}{chapter}{如何使用这本讲义}

这本讲义面向第一次系统学习信息论、同时希望把知识用于随机算法、数据流 Sketch、AI4Math 与 AI4TCS 的读者。它不要求测度论，默认随机变量取有限字母表；所有对数若无特别说明均以 2 为底，因此信息量单位为 bit。

建议每章按四步学习：
\begin{enumerate}
  \item 先写出本章要回答的问题，不先背公式；
  \item 手算正文中的例子，确认公式中的随机变量和分布；
  \item 阅读证明骨架，标出每一步使用的假设；
  \item 完成章末验收，再进入下一章。
\end{enumerate}

讲义中的“AI4Math / AI4TCS 连接”主要用于提出可检验问题。除非明确给出定理和假设，否则这些连接不应被当作已有理论结论。

\tableofcontents
\clearpage

\chapter{概率、信息量与熵}

\section{这一章回答什么}

信息论首先要回答一个看似简单的问题：一个随机结果出现时，我们“获得了多少信息”？回答这个问题前，必须先把随机对象、分布和平均方式写清楚。本章集中处理单个离散随机变量，目标是建立后续所有公式的共同语言。

\section{有限概率空间与随机变量}

设样本空间为有限集合 $\Omega$，概率函数 $\Prb$ 满足
\[
\Prb(\omega)\ge 0,
\qquad
\sum_{\omega\in\Omega}\Prb(\omega)=1.
\]
随机变量 $X:\Omega\to\mathcal X$ 把基本结果映射到字母表 $\mathcal X$。它的概率质量函数为
\[
p_X(x)=\Prb[X=x].
\]

对任意函数 $f$，期望定义为
\[
\E[f(X)]=\sum_{x\in\mathcal X}p_X(x)f(x).
\]
方差衡量围绕均值的平方波动：
\[
\Var(X)=\E[(X-\E X)^2]=\E[X^2]-(\E X)^2.
\]

\begin{example}[Bernoulli 随机变量]
若 $X\sim\mathrm{Bernoulli}(p)$，则
\[
\Prb[X=1]=p,\qquad \Prb[X=0]=1-p,
\]
并且 $\E X=p$、$\Var(X)=p(1-p)$。同一个取值 $X=1$ 在不同 $p$ 下有不同的意外程度；信息量因此必须依赖分布，而不能只依赖符号本身。
\end{example}

\section{从意外程度到自信息}

如果事件发生概率为 $p$，希望它的信息量 $i(p)$ 满足：
\begin{enumerate}
  \item 概率越小，信息量越大；
  \item 必然事件的信息量为 $0$；
  \item 独立事件同时发生时，信息量相加，即 $i(pq)=i(p)+i(q)$；
  \item $i(p)$ 随 $p$ 连续变化。
\end{enumerate}
这些条件把函数确定为对数形式。

\begin{definition}[自信息]
事件 $X=x$ 的自信息定义为
\[
\imath_X(x)=-\log_2 p_X(x).
\]
当 $p_X(x)=0$ 时，该取值不会被实际观测；求和中采用约定 $0\log 0=0$。
\end{definition}

若 $p_X(x)=1/2$，观测它得到 $1$ bit；若概率为 $1/8$，得到 $3$ bits。这里的 bit 是信息单位，不是变量在计算机中实际占用的位数。

\begin{warningbox}
自信息描述“在给定分布模型下，一个结果有多意外”。它不等于语义重要性。一个极小概率的乱码可能有很高自信息，但对研究问题没有任何价值。
\end{warningbox}

\section{熵：平均自信息}

\begin{definition}[Shannon 熵]
离散随机变量 $X$ 的熵为
\[
H(X)=\E[\imath_X(X)]
=-\sum_{x\in\mathcal X}p_X(x)\log_2 p_X(x).
\]
\end{definition}

熵是分布 $p_X$ 的函数。它对符号名字不敏感，只关心各概率的排列。

\subsection{二元熵函数}

对 $X\sim\mathrm{Bernoulli}(p)$，定义
\[
h_2(p)=-p\log_2 p-(1-p)\log_2(1-p).
\]
它满足 $h_2(0)=h_2(1)=0$、$h_2(1/2)=1$，并关于 $p=1/2$ 对称。

对 $0<p<1$ 求导：
\[
h_2'(p)=\log_2\frac{1-p}{p},
\qquad
h_2''(p)=-\frac{1}{\ln 2}\left(\frac1p+\frac1{1-p}\right)<0.
\]
因此二元熵严格凹，并在 $p=1/2$ 取得最大值。

\begin{example}[偏斜硬币]
若 $p=1/4$，则
\[
h_2(1/4)
=-\frac14\log_2\frac14-\frac34\log_2\frac34
\approx 0.811\bits.
\]
它小于公平硬币的 $1$ bit，因为偏斜硬币更容易预测。
\end{example}

\section{熵的基本边界}

\begin{theorem}[有限字母表上的熵界]
若 $|\mathcal X|=m$，则
\[
0\le H(X)\le \log_2 m.
\]
左侧等号当且仅当 $X$ 几乎必然为某个固定值；右侧等号当且仅当 $X$ 在 $\mathcal X$ 上均匀分布。
\end{theorem}

\begin{proof}[证明骨架]
每一项 $-p(x)\log p(x)$ 非负，所以 $H(X)\ge 0$。上界可以由下一章的 KL 非负性得到。令 $U$ 为 $\mathcal X$ 上的均匀分布，则
\[
D_{\mathrm{KL}}(P_X\|U)
=\sum_x p(x)\log_2\frac{p(x)}{1/m}
=\log_2 m-H(X)\ge 0.
\]
等号条件来自 KL 散度为零当且仅当两个分布相同。
\end{proof}

\begin{intuitionbox}
熵上界说的是：若只有 $m$ 种可能结果，平均不确定性不会超过区分 $m$ 个等概率结果所需的 $\log_2 m$ bits。它不是说任何单个符号都能用 $\log_2 m$ 个比特无损编码；编码是否为整数长度、是否允许分块，还要在第三章讨论。
\end{intuitionbox}

\section{独立重复与可加性}

若 $X_1,\dots,X_n$ 相互独立，则联合分布分解为
\[
p(x_1,\dots,x_n)=\prod_{i=1}^n p_i(x_i).
\]
于是
\begin{align*}
H(X_1,\dots,X_n)
&=-\E\log_2\prod_{i=1}^n p_i(X_i)\\
&=\sum_{i=1}^n H(X_i).
\end{align*}
对独立同分布样本 $X^n=(X_1,\dots,X_n)$，有 $H(X^n)=nH(X)$。

\begin{warningbox}
没有独立性时，联合熵一般不等于边缘熵之和。正确关系是下一章的链式法则
$H(X,Y)=H(X)+H(Y|X)$。
\end{warningbox}

\section{一个完整数值例子}

设字母表为 $\{a,b,c,d\}$，分布为
\[
p=(1/2,1/4,1/8,1/8).
\]
四个符号的自信息分别为 $1,2,3,3$ bits，因此
\[
H(X)=\frac12\cdot1+\frac14\cdot2+\frac18\cdot3+\frac18\cdot3
=1.75\bits.
\]
这个例子随后会对应一棵 Huffman 树，其码长正好为 $1,2,3,3$，平均码长等于熵。一般分布不一定恰好达到熵，只能满足相应上下界。

\section{小结与验收}

\begin{checkpointbox}
在不看讲义的情况下完成：
\begin{enumerate}
  \item 解释自信息和熵分别是“单次量”还是“平均量”；
  \item 推导 $h_2'(p)$ 并说明公平硬币为何熵最大；
  \item 手算分布 $(1/2,1/4,1/8,1/8)$ 的熵；
  \item 说明“熵为 1 bit”为什么不等于每个观测在文件中恰好占 1 bit。
\end{enumerate}
\end{checkpointbox}

\chapter{联合熵、互信息与 KL 散度}

上一章只处理一个随机变量。本章研究两个变量之间的信息关系：知道 $Y$ 后，$X$ 还剩多少不确定性？两个变量共享多少统计信息？两个分布相差多少？

\section{联合分布、边缘分布与条件分布}

联合概率质量函数为 $p_{XY}(x,y)$。边缘分布由求和得到：
\[
p_X(x)=\sum_y p_{XY}(x,y),
\qquad
p_Y(y)=\sum_x p_{XY}(x,y).
\]
当 $p_Y(y)>0$ 时，条件分布为
\[
p_{X|Y}(x|y)=\frac{p_{XY}(x,y)}{p_Y(y)}.
\]
独立性等价于 $p_{XY}(x,y)=p_X(x)p_Y(y)$ 对所有 $(x,y)$ 成立。

\begin{definition}[联合熵与条件熵]
\begin{align*}
H(X,Y)&=-\sum_{x,y}p(x,y)\log_2p(x,y),\\
H(X|Y)&=\sum_y p(y)H(X|Y=y)\\
&=-\sum_{x,y}p(x,y)\log_2p(x|y).
\end{align*}
\end{definition}

\section{熵的链式法则}

\begin{theorem}[链式法则]
\[
H(X,Y)=H(Y)+H(X|Y)=H(X)+H(Y|X).
\]
更一般地，
\[
H(X_1,\dots,X_n)=\sum_{i=1}^n H(X_i|X_1,\dots,X_{i-1}).
\]
\end{theorem}

\begin{proof}
利用 $p(x,y)=p(y)p(x|y)$，有
\begin{align*}
H(X,Y)
&=-\sum_{x,y}p(x,y)\log_2[p(y)p(x|y)]\\
&=-\sum_{x,y}p(x,y)\log_2p(y)
  -\sum_{x,y}p(x,y)\log_2p(x|y)\\
&=H(Y)+H(X|Y).
\end{align*}
交换 $X,Y$ 得到另一种分解。多变量版本反复应用二变量链式法则即可。
\end{proof}

\begin{intuitionbox}
联合不确定性可以按顺序支付：先描述 $Y$，再在已知 $Y$ 的条件下描述 $X$。链式法则是编码、典型集和通信协议中最常用的“信息记账规则”。
\end{intuitionbox}

\section{KL 散度与 Gibbs 不等式}

\begin{definition}[KL 散度]
对有限集合上的分布 $P,Q$，若 $P(x)>0$ 时均有 $Q(x)>0$，定义
\[
D_{\mathrm{KL}}(P\|Q)
=\sum_x P(x)\log_2\frac{P(x)}{Q(x)}.
\]
若存在 $P(x)>0$ 但 $Q(x)=0$，则散度为 $+\infty$。
\end{definition}

\begin{theorem}[Gibbs 不等式]
\[
D_{\mathrm{KL}}(P\|Q)\ge 0,
\]
等号当且仅当 $P=Q$（在 $P$ 的支持上）。
\end{theorem}

\begin{proof}
使用自然对数不等式 $\ln t\le t-1$：
\begin{align*}
D_{\mathrm{KL}}(P\|Q)
&=-\frac1{\ln2}\sum_xP(x)\ln\frac{Q(x)}{P(x)}\\
&\ge -\frac1{\ln2}\sum_xP(x)
\left(\frac{Q(x)}{P(x)}-1\right)\\
&=-\frac1{\ln2}\left(\sum_xQ(x)-\sum_xP(x)\right)=0.
\end{align*}
等号要求 $Q(x)/P(x)=1$ 对所有 $P(x)>0$ 成立。
\end{proof}

\subsection{交叉熵}

交叉熵定义为
\[
H(P,Q)=-\sum_xP(x)\log_2Q(x).
\]
它满足
\[
H(P,Q)=H(P)+D_{\mathrm{KL}}(P\|Q).
\]
当真实分布固定时，最小化交叉熵等价于最小化 $D_{\mathrm{KL}}(P\|Q)$。这解释了分类模型中的交叉熵目标，但不意味着交叉熵下降自动带来推理、证明或算法能力。

\begin{warningbox}
KL 散度不是距离：通常 $D(P\|Q)\ne D(Q\|P)$，也不满足三角不等式。书写时不能省略方向。
\end{warningbox}

\begin{example}[KL 的不对称性]
取 $P=(1/2,1/2)$、$Q=(3/4,1/4)$，则
\begin{align*}
D(P\|Q)&=\tfrac12\log_2\tfrac23+\tfrac12\log_2 2\approx0.2075,\\
D(Q\|P)&=\tfrac34\log_2\tfrac32+\tfrac14\log_2\tfrac12\approx0.1887.
\end{align*}
两者不同。
\end{example}

\section{互信息}

\begin{definition}[互信息]
\[
I(X;Y)=D_{\mathrm{KL}}(P_{XY}\|P_XP_Y)
=\sum_{x,y}p(x,y)\log_2\frac{p(x,y)}{p(x)p(y)}.
\]
\end{definition}

由 KL 非负性，$I(X;Y)\ge0$；等号当且仅当 $X,Y$ 独立。

把联合分布分解为 $p(x,y)=p(y)p(x|y)$，可得四个等价表达：
\begin{align*}
I(X;Y)
&=H(X)-H(X|Y)\\
&=H(Y)-H(Y|X)\\
&=H(X)+H(Y)-H(X,Y)\\
&=D_{\mathrm{KL}}(P_{XY}\|P_XP_Y).
\end{align*}

因此，互信息既可以解释为“观察 $Y$ 后 $X$ 的不确定性平均减少多少”，也可以解释为“联合分布偏离独立分布多远”。

\begin{example}[带噪复制]
设 $X\sim\mathrm{Bernoulli}(1/2)$，$N\sim\mathrm{Bernoulli}(q)$ 且独立，令 $Y=X\oplus N$。则 $Y$ 仍均匀，
\[
H(Y)=1,\qquad H(Y|X)=H(N)=h_2(q),
\]
所以
\[
I(X;Y)=1-h_2(q).
\]
当 $q=0$ 时完全保留 1 bit；当 $q=1/2$ 时输出与输入独立，互信息为 0。
\end{example}

\section{条件互信息与链式法则}

\begin{definition}[条件互信息]
\[
I(X;Y|Z)=H(X|Z)-H(X|Y,Z).
\]
\end{definition}
它是对每个 $Z=z$ 的互信息取平均，因此非负。互信息链式法则为
\[
I(X;Y,Z)=I(X;Z)+I(X;Y|Z).
\]
进一步，
\[
I(X;Y_1,\dots,Y_n)
=\sum_{i=1}^n I(X;Y_i|Y_1,\dots,Y_{i-1}).
\]

\section{条件化减少熵与数据处理不等式}

由 $I(X;Y)=H(X)-H(X|Y)\ge0$，得到
\[
H(X|Y)\le H(X).
\]
这是一条平均意义的结论；对某个特定 $Y=y$，条件熵可以高于原熵。

\begin{theorem}[数据处理不等式]
若 $X\to Y\to Z$ 构成 Markov 链，即给定 $Y$ 后 $Z$ 与 $X$ 条件独立，则
\[
I(X;Z)\le I(X;Y).
\]
\end{theorem}

\begin{proof}[证明骨架]
一方面，由链式法则
\[
I(X;Y,Z)=I(X;Y)+I(X;Z|Y)=I(X;Y),
\]
其中 Markov 条件给出 $I(X;Z|Y)=0$。另一方面，
\[
I(X;Y,Z)=I(X;Z)+I(X;Y|Z)\ge I(X;Z).
\]
合并即得结论。
\end{proof}

\begin{intuitionbox}
对数据做后处理不能凭空增加它与原始变量之间的统计信息。后处理可以重新组织、压缩或突出某些信息，但若只观察 $Z=f(Y)$，关于 $X$ 的总互信息不会超过观察完整 $Y$。
\end{intuitionbox}

\section{一个联合分布例子}

令联合分布为
\[
\begin{array}{c|cc}
 &Y=0&Y=1\\\hline
X=0&3/8&1/8\\
X=1&1/8&3/8
\end{array}
\]
两个边缘分布都均匀，所以 $H(X)=H(Y)=1$。给定 $Y$ 后，$X=Y$ 的概率为 $3/4$，故
\[
H(X|Y)=h_2(1/4)\approx0.811,
\qquad
I(X;Y)\approx0.189\bits.
\]
这个数值等于一个交叉概率 $q=1/4$ 的二元对称信道所保留的信息。

\section{小结与验收}

\begin{checkpointbox}
\begin{enumerate}
  \item 从 $p(x,y)=p(y)p(x|y)$ 推导熵的链式法则；
  \item 说明 KL 散度非负的关键不等式和等号条件；
  \item 写出互信息的四种等价表达，并解释每一种视角；
  \item 对上面的 $2\times2$ 联合分布手算 $H(X|Y)$ 和 $I(X;Y)$；
  \item 用一句话说明数据处理不等式依赖的 Markov 假设。
\end{enumerate}
\end{checkpointbox}

\chapter{无失真信源编码：从前缀码到 Huffman}

\section{编码问题的完整定义}

设信源符号 $X$ 取值于有限字母表 $\mathcal X=\{x_1,\dots,x_m\}$，概率为 $p_i=\Prb[X=x_i]$。二元变长码是映射
\[
C:\mathcal X\to\{0,1\}^*,
\]
其中 $\{0,1\}^*$ 表示所有有限二进制串。记码字长度为 $\ell_i=|C(x_i)|$，平均码长为
\[
L(C)=\sum_{i=1}^m p_i\ell_i.
\]

几个常被混用的概念必须分开：
\begin{description}
  \item[非奇异（nonsingular）] 不同源符号映到不同码字；
  \item[唯一可译（uniquely decodable）] 任意有限源序列连接后的比特串只能被解成唯一源序列；
  \item[前缀码（prefix-free）] 没有码字是另一码字的前缀；
  \item[瞬时码（instantaneous）] 从左到右读到一个码字结尾时即可立即输出符号，不必等待后续比特。
\end{description}
二元前缀码与瞬时码是同一类。前缀码一定唯一可译，但唯一可译码未必前缀无关；非奇异也远弱于唯一可译。

\begin{example}[非奇异但不可唯一译]
取 $C(a)=0,C(b)=01,C(c)=10$。码字互不相同，但串 $010$ 既可译为 $ac$，也可译为 $ba$，所以不是唯一可译码。
\end{example}

\section{前缀码与二叉树}

把比特 0/1 看成从根节点向左/右走一步。每个码字对应一条根到节点的路径。前缀码要求所有码字都位于叶节点：如果某码字对应内部节点，它就会成为更长码字的前缀。

\begin{center}
\begin{tikzpicture}[
  level distance=10mm,
  level 1/.style={sibling distance=42mm},
  level 2/.style={sibling distance=24mm},
  level 3/.style={sibling distance=13mm},
  every node/.style={font=\small}
]
\node[circle,draw] {}
 child {node[circle,draw] {$a$} edge from parent node[left] {0}}
 child {node[circle,draw] {}
   child {node[circle,draw] {$b$} edge from parent node[left] {0}}
   child {node[circle,draw] {}
     child {node[circle,draw] {$c$} edge from parent node[left] {0}}
     child {node[circle,draw] {$d$} edge from parent node[right] {1}}
     edge from parent node[right] {1}}
   edge from parent node[right] {1}};
\end{tikzpicture}
\end{center}
这棵树对应码字 $0,10,110,111$，码长为 $1,2,3,3$。

\section{Kraft--McMillan 不等式}

\begin{theorem}[Kraft 不等式]
若存在码长为 $\ell_1,\dots,\ell_m$ 的二元前缀码，则
\[
\sum_{i=1}^m 2^{-\ell_i}\le1.
\]
反过来，只要一组正整数码长满足该不等式，就存在具有这些码长的二元前缀码。
\end{theorem}

\begin{proof}[树的证明]
取 $L=\max_i\ell_i$。深度为 $\ell_i$ 的叶节点覆盖深度 $L$ 上的 $2^{L-\ell_i}$ 个后代。前缀无关保证这些后代集合互不相交，因此
\[
\sum_i2^{L-\ell_i}\le2^L.
\]
两边除以 $2^L$ 得必要性。充分性可按码长从短到长贪心分配尚未被占用的叶节点；Kraft 和不超过 1 保证每一步都有位置。
\end{proof}

对唯一可译码，更一般的 McMillan 不等式仍给出同样的必要条件。因此，在最小化平均码长时，限制到前缀码不会损失最优值。

\begin{intuitionbox}
$2^{-\ell_i}$ 可以理解为码字在无限二叉树边界上占据的“份额”。短码字占据更大的子树，所以不能给太多符号同时分配短码字。
\end{intuitionbox}

\section{熵与平均码长的基本界}

\begin{theorem}[无失真一符号编码界]
对任意唯一可译二元码，
\[
L\ge H(X).
\]
并且存在前缀码满足
\[
H(X)\le L < H(X)+1.
\]
\end{theorem}

\begin{proof}[下界]
令 $K=\sum_i2^{-\ell_i}\le1$，定义分布
\[
q_i=\frac{2^{-\ell_i}}{K}.
\]
直接计算可得
\[
D_{\mathrm{KL}}(P\|Q)=L-H(X)+\log_2K.
\]
因此
\[
L-H(X)=D_{\mathrm{KL}}(P\|Q)-\log_2K\ge0.
\]
\end{proof}

\begin{proof}[上界]
取 Shannon 码长
\[
\ell_i=\left\lceil-\log_2p_i\right\rceil.
\]
则 $2^{-\ell_i}\le p_i$，所以 Kraft 和不超过 1，存在相应前缀码。同时
\[
\ell_i<-\log_2p_i+1,
\]
对 $p_i$ 加权求和即得 $L<H(X)+1$。
\end{proof}

\begin{warningbox}
$L\ge H(X)$ 是平均码长的结论，不是说每个符号码长都至少为其自信息，也不是说任意一个文件都一定能压缩。整数码长造成的一符号冗余可以通过分块编码摊薄。
\end{warningbox}

\section{Huffman 算法}

Huffman 算法在所有二元前缀码中最小化平均码长：
\begin{enumerate}
  \item 选择概率最小的两个符号；
  \item 把它们合并为概率之和的新节点；
  \item 在缩小后的字母表上重复；
  \item 逆序展开合并树，左右边分别标 0/1。
\end{enumerate}

\begin{example}[完整 Huffman 计算]
仍取 $p=(1/2,1/4,1/8,1/8)$：
\[
1/8+1/8\to1/4,
\qquad
1/4+1/4\to1/2,
\qquad
1/2+1/2\to1.
\]
得到码长 $1,2,3,3$，平均码长
\[
L=\frac12+\frac12+\frac38+\frac38=1.75\bits=H(X).
\]
这里恰好等于熵，是因为所有概率都是 $2$ 的负整数次幂。
\end{example}

\subsection{为什么贪心是最优的}

最优性证明包含两个关键事实：
\begin{enumerate}
  \item 存在一棵最优前缀树，使概率最小的两个符号位于最深层且互为兄弟；这是通过交换叶节点而不增大平均码长得到的；
  \item 把这两个兄弟合并为一个伪符号后，原问题的最优树对应缩小问题的最优树，否则可替换为更优子树并产生矛盾。
\end{enumerate}
这正是贪心选择性质与最优子结构。

若使用最小堆，Huffman 建树时间为 $O(m\log m)$。编码和解码时间还取决于消息长度和树/码表表示。

\section{分块编码与渐近信源编码定理}

对独立同分布信源 $X_1,X_2,\dots$，把长度 $n$ 的块 $X^n$ 当作一个大符号。直接应用上一节的界：
\[
H(X^n)\le L_n<H(X^n)+1.
\]
因为独立同分布给出 $H(X^n)=nH(X)$，故每个源符号的平均码长满足
\[
H(X)\le\frac{L_n}{n}<H(X)+\frac1n.
\]
这说明整数码长造成的每符号冗余可随块长趋于无穷而消失。

\subsection{典型集与 AEP}

定义弱典型集
\[
A_\epsilon^{(n)}=
\left\{x^n:
\left|-\frac1n\log_2p(x^n)-H(X)\right|\le\epsilon
\right\}.
\]
渐近等概率性质（AEP）说明：对任意 $\epsilon>0$，
\[
\Prb[X^n\in A_\epsilon^{(n)}]\to1.
\]
典型序列的概率约为 $2^{-nH(X)}$，典型集大小约为 $2^{nH(X)}$。更精确地，当 $n$ 足够大时，
\[
(1-\epsilon)2^{n(H-\epsilon)}
\le |A_\epsilon^{(n)}|
\le 2^{n(H+\epsilon)}.
\]

\begin{theorem}[几乎无失真固定长度信源编码]
对独立同分布有限字母表信源：
\begin{itemize}
  \item 若码率 $R>H(X)$，存在长度 $n$ 的固定长度块码，使错误概率趋于 0；
  \item 若 $R<H(X)$，错误概率不能趋于 0。
\end{itemize}
\end{theorem}

\begin{proof}[可达性骨架]
给每个典型序列分配唯一索引，非典型序列统一报错。典型集大小至多
\[
2^{n(H+\epsilon)}.
\]
所以当 $R>H+\epsilon$ 时有足够的 $2^{nR}$ 个索引。错误只在非典型事件上发生，其概率趋于 0。
\end{proof}

\begin{proof}[逆定理直觉]
码率 $R$ 的固定长度码最多区分 $2^{nR}$ 个序列，而高概率质量分散在约 $2^{nH}$ 个典型序列上。若 $R<H$，可唯一表示的典型序列只占指数级变小的一部分，无法让错误概率消失。
\end{proof}

\section{三类结论不要混在一起}

\begin{center}
\begin{tabularx}{\textwidth}{lXXX}
\toprule
情形 & 是否允许错误 & 主要对象 & 结论\\
\midrule
一符号变长码 & 不允许 & 平均整数码长 & $H\le L<H+1$\\
分块变长码 & 不允许 & 每符号平均码长 & $H\le L_n/n<H+1/n$\\
固定长度块码 & 允许小错误 & 码率与错误概率 & $R>H$ 可使错误趋零\\
\bottomrule
\end{tabularx}
\end{center}

\section{小结与验收}

\begin{checkpointbox}
\begin{enumerate}
  \item 给出“非奇异、唯一可译、前缀码”的包含关系和一个反例；
  \item 用二叉树证明 Kraft 不等式；
  \item 对概率 $(0.4,0.3,0.2,0.1)$ 手工运行 Huffman 算法并计算平均码长；
  \item 解释 $H\le L<H+1$ 中的 1 从哪里来，分块为何能摊薄它；
  \item 区分无错变长编码定理与允许小错误的固定长度信源编码定理。
\end{enumerate}
\end{checkpointbox}

\chapter{离散信道、互信息与容量}

信源编码问“如何消除冗余”；信道编码问“如何在噪声中可靠通信”。两者使用同一个熵与互信息语言，但优化对象和允许的错误不同。

\section{离散无记忆信道}

\begin{definition}[离散无记忆信道]
输入字母表为 $\mathcal X$、输出字母表为 $\mathcal Y$。信道由转移概率
\[
W(y|x)=\Prb[Y=y|X=x]
\]
描述，并满足对每个 $x$ 有 $\sum_yW(y|x)=1$。无记忆表示长度 $n$ 使用时
\[
W^n(y^n|x^n)=\prod_{i=1}^nW(y_i|x_i).
\]
\end{definition}

给定输入分布 $p_X$ 后，联合分布为
\[
p_{XY}(x,y)=p_X(x)W(y|x),
\]
输出分布为 $p_Y(y)=\sum_xp_X(x)W(y|x)$。互信息 $I(X;Y)$ 衡量一次信道使用平均保留了多少关于输入的信息。

\section{码、码率与错误概率}

一个 $(n,M)$ 信道码包含：
\begin{itemize}
  \item 消息 $J$ 在 $\{1,\dots,M\}$ 中取值；
  \item 编码器 $f:J\mapsto X^n$；
  \item 信道产生 $Y^n\sim W^n(\cdot|X^n)$；
  \item 解码器 $g:Y^n\mapsto\hat J$。
\end{itemize}
码率为
\[
R=\frac1n\log_2M\quad\text{bits/channel use}.
\]
平均错误概率为 $P_e=\Prb[\hat J\ne J]$。最大错误概率则对最坏消息取最大值；两种口径不能在结果表中混用。

\section{二元对称信道 BSC}

BSC$(q)$ 以概率 $q$ 翻转输入比特：
\[
Y=X\oplus Z,\qquad Z\sim\mathrm{Bernoulli}(q),
\]
且 $Z$ 与 $X$ 独立。

\begin{center}
\begin{tikzpicture}[node distance=34mm,>=Latex]
\node (x0) {$X=0$};
\node[below=14mm of x0] (x1) {$X=1$};
\node[right=of x0] (y0) {$Y=0$};
\node[below=14mm of y0] (y1) {$Y=1$};
\draw[->] (x0) -- node[above] {$1-q$} (y0);
\draw[->] (x0) -- node[near start,below,sloped] {$q$} (y1);
\draw[->] (x1) -- node[near start,above,sloped] {$q$} (y0);
\draw[->] (x1) -- node[below] {$1-q$} (y1);
\end{tikzpicture}
\end{center}

因为给定 $X$ 后 $Y$ 的不确定性只来自噪声，
\[
H(Y|X)=H(Z)=h_2(q).
\]
所以
\[
I(X;Y)=H(Y)-h_2(q).
\]

\section{信道容量}

\begin{definition}[容量]
离散无记忆信道的容量定义为
\[
C=\max_{p_X}I(X;Y).
\]
最大化变量是输入分布，而信道 $W(y|x)$ 固定。
\end{definition}

\begin{theorem}[BSC 容量]
对 $0\le q\le1/2$，
\[
C_{\mathrm{BSC}(q)}=1-h_2(q).
\]
均匀输入 $X\sim\mathrm{Bernoulli}(1/2)$ 达到容量。
\end{theorem}

\begin{proof}
由 $I(X;Y)=H(Y)-h_2(q)$ 且二元变量熵至多为 1，
\[
I(X;Y)\le1-h_2(q).
\]
当输入均匀时，BSC 的对称性使输出也均匀，故 $H(Y)=1$，达到上界。
\end{proof}

边界情况给出直观检查：$q=0$ 时无噪声，容量为 1；$q=1/2$ 时输出独立于输入，容量为 0。若 $q>1/2$，接收端可以统一翻转输出，把交叉概率降为 $1-q$。

\section{二元擦除信道 BEC}

BEC$(\varepsilon)$ 以概率 $1-\varepsilon$ 正确输出输入比特，以概率 $\varepsilon$ 输出擦除符号 $?$。接收端知道哪些位置被擦除。

给定均匀输入，未擦除时获得 1 bit，擦除时获得 0 bit，因此
\[
C_{\mathrm{BEC}(\varepsilon)}=1-\varepsilon.
\]
BEC 与 BSC 的区别在于：擦除位置可见，而翻转位置不可见。相同数值的 $q$ 与 $\varepsilon$ 不代表相同难度。

\section{容量究竟说明什么}

容量是一个渐近阈值：
\begin{itemize}
  \item 当 $R<C$，存在一列块长增长的码，使错误概率趋于 0；
  \item 当 $R>C$，不可能让错误概率趋于 0。
\end{itemize}
容量定义本身只给出单字母优化问题。把它解释成可靠通信阈值，需要下一章的信道编码定理。

\begin{warningbox}
容量不等于某个具体码在有限块长上的吞吐量，也不包含编码/解码计算成本。随机编码证明“存在好码”，不代表已经给出可高效实现的码。
\end{warningbox}

\section{带输入成本约束的容量}

若输入符号 $x$ 有成本 $c(x)$，并要求 $\E[c(X)]\le\Gamma$，容量变为
\[
C(\Gamma)=\max_{p_X:\,\E c(X)\le\Gamma}I(X;Y).
\]
这提醒我们：优化信息量时必须同时写清资源约束。AI 系统中的 token、查询、编译或 verifier 调用也有成本，但只有在建立了精确随机模型后，才能合法地套用信息论结论。

\section{小结与验收}

\begin{checkpointbox}
\begin{enumerate}
  \item 写出 DMC、$(n,M)$ 码、码率和平均错误概率的定义；
  \item 从 $I(X;Y)=H(Y)-H(Y|X)$ 推导 BSC 容量；
  \item 比较 BSC 与 BEC：接收端分别知道什么；
  \item 解释容量是渐近存在性阈值，而不是有限长度代码的直接性能；
  \item 给一个需要成本约束的现实通信或算法例子。
\end{enumerate}
\end{checkpointbox}

\chapter{随机编码与可靠通信}

本章解释信道编码定理的两半：为什么低于容量的速率存在可靠码，以及为什么高于容量不可能可靠通信。核心工具分别是随机码本与典型性、Fano 不等式与数据处理。

\section{联合典型性}

对输入分布 $p_X$ 和信道 $W(y|x)$，联合分布为 $p_{XY}(x,y)=p_X(x)W(y|x)$。长度 $n$ 的序列对 $(x^n,y^n)$ 若其经验信息量接近相应熵，称为联合典型。

本章只使用联合典型性的三个性质：
\begin{enumerate}
  \item 若 $X^n\sim p_X^n$ 且 $Y^n\sim W^n(\cdot|X^n)$，则 $(X^n,Y^n)$ 以趋于 1 的概率联合典型；
  \item 若 $\widetilde X^n\sim p_X^n$ 与 $Y^n\sim p_Y^n$ 独立，则
  \[
  \Prb[(\widetilde X^n,Y^n)\text{ 联合典型}]
  \lesssim 2^{-n(I(X;Y)-\delta(\epsilon))};
  \]
  \item 典型序列的概率和集合大小都具有指数阶控制。
\end{enumerate}
第二条最关键：一个与输出无关的错误码字，伪装成“与输出匹配”的概率约为 $2^{-nI(X;Y)}$。

\section{随机码本与典型译码}

固定输入分布 $p_X$ 和码率 $R$，令 $M=2^{nR}$。随机生成码本：对每个消息 $m\in\{1,\dots,M\}$，独立采样
\[
X^n(m)\sim\prod_{i=1}^np_X(x_i).
\]
发送消息 $m$ 时把 $X^n(m)$ 输入信道。解码器观察 $Y^n$ 后，寻找唯一满足 $(X^n(\hat m),Y^n)$ 联合典型的码字；若不存在或不唯一则报错。

由于码本生成和信道都对消息编号对称，可以假设发送消息 1。定义错误事件：
\begin{align*}
E_0&=\{(X^n(1),Y^n)\text{ 不联合典型}\},\\
E_m&=\{(X^n(m),Y^n)\text{ 联合典型}\},\quad m\ne1.
\end{align*}
整体错误包含于
\[
E_0\cup\bigcup_{m=2}^ME_m.
\]

\section{Union bound 给出的速率条件}

由联合典型性，$\Prb(E_0)\to0$。对于每个错误码字，由于 $X^n(m)$ 与实际输出 $Y^n$ 独立，
\[
\Prb(E_m)\lesssim2^{-n(I(X;Y)-\delta)}.
\]
使用 union bound：
\begin{align*}
\Prb[\text{error}]
&\le\Prb(E_0)+\sum_{m=2}^M\Prb(E_m)\\
&\lesssim\Prb(E_0)+2^{nR}2^{-n(I(X;Y)-\delta)}\\
&=\Prb(E_0)+2^{-n(I(X;Y)-R-\delta)}.
\end{align*}
只要
\[
R<I(X;Y)-\delta,
\]
平均错误概率就趋于 0。再对输入分布优化，得到所有 $R<C=\max_{p_X}I(X;Y)$ 都可达。

\begin{theorem}[信道编码定理：可达性]
对离散无记忆信道，任意 $R<C$ 都存在一列 $(n,2^{nR})$ 码，使平均错误概率随 $n\to\infty$ 趋于 0。
\end{theorem}

\section{为什么随机证明能推出确定码存在}

上面的错误概率同时对随机码本、随机消息和随机信道噪声取平均。若所有确定码本的错误概率都大于某个数，随机选择后的平均也会大于该数。反过来，随机码本平均错误概率很小，必然至少存在一个确定码本不超过这个平均值。

\begin{intuitionbox}
随机编码是概率方法：随机性主要服务于存在性证明。证明结束后，可以固定一个表现良好的码本。它没有自动告诉我们如何高效找到、存储或解码这个码本。
\end{intuitionbox}

若要从平均错误转成最大错误，还需删除错误概率最差的一部分码字（expurgation）或使用更细分析。论文中必须明确采用哪种错误口径。

\section{Fano 不等式}

\begin{theorem}[Fano 不等式]
设消息 $J$ 取 $M$ 个值，估计为 $\hat J$，错误概率 $P_e=\Prb[\hat J\ne J]$，则
\[
H(J|\hat J)\le h_2(P_e)+P_e\log_2(M-1)
\le1+P_e\log_2M.
\]
\end{theorem}

\begin{proof}[证明骨架]
令错误指示变量 $E=\ind\{J\ne\hat J\}$。一方面，
\[
H(E,J|\hat J)=H(J|\hat J)+H(E|J,\hat J)=H(J|\hat J),
\]
因为 $E$ 由 $(J,\hat J)$ 决定。另一方面，
\[
H(E,J|\hat J)=H(E|\hat J)+H(J|E,\hat J).
\]
第一项至多 $h_2(P_e)$；当 $E=0$ 时 $J$ 已知，当 $E=1$ 时至多剩 $M-1$ 种可能，所以第二项至多 $P_e\log_2(M-1)$。
\end{proof}

Fano 不等式把“低错误率”转换为“小条件熵”：如果几乎总能从观察中恢复消息，那么观察后关于消息剩余的不确定性必须很小。

\section{信道编码逆定理}

设消息 $J$ 均匀分布，因此 $H(J)=\log_2M=nR$。编码和信道形成 Markov 链
\[
J\to X^n\to Y^n\to\hat J.
\]
由熵分解与 Fano 不等式：
\begin{align*}
nR
&=I(J;Y^n)+H(J|Y^n)\\
&\le I(X^n;Y^n)+1+P_enR.
\end{align*}
对无记忆信道，链式法则和容量定义给出
\[
I(X^n;Y^n)\le\sum_{i=1}^nI(X_i;Y_i)\le nC.
\]
因此
\[
R\le C+\frac1n+P_eR.
\]
若 $P_e\to0$ 且 $n\to\infty$，必有 $R\le C$。

\begin{theorem}[信道编码定理：逆定理]
若存在一列码使错误概率趋于 0，则其渐近码率不超过容量 $C$。
\end{theorem}

\section{存在性、构造性与有限块长}

信道编码定理是渐近结果，不能直接回答有限 $n$ 的工程问题。实际还要考虑：
\begin{itemize}
  \item 码本构造和存储成本；
  \item 编码、解码复杂度与延迟；
  \item 有限块长下的错误概率；
  \item 平均错误与最大错误；
  \item 信道模型失配和输入成本约束。
\end{itemize}
现代编码理论中的线性码、LDPC、polar code 等，正是在可靠性、码率和计算效率之间寻找可构造方案。

\begin{warningbox}
“随机搜索找到一个在测试集上正确的程序”和随机编码定理不是同一件事。后者有明确的概率空间、错误事件、指数规模码本和渐近证明；前者若没有隐藏测试、分布假设与理论分析，只是经验结果。
\end{warningbox}

\section{小结与验收}

\begin{checkpointbox}
\begin{enumerate}
  \item 画出随机码本、信道、典型译码的完整流程；
  \item 写出 $E_0$ 与 $E_m$，说明它们为何具有不同概率界；
  \item 从 union bound 推出条件 $R<I(X;Y)$；
  \item 复述 Fano 不等式证明中错误指示变量的作用；
  \item 解释“随机码本平均表现好”为什么能推出至少一个确定码本存在；
  \item 列出信道编码定理没有解决的三个有限块长问题。
\end{enumerate}
\end{checkpointbox}

\chapter{信息论与 Sketch、通信复杂度和 AI4TCS}

本章不再引入完整新定理，而是把前五章的对象映射到数据流和可验证算法研究中。每个连接都要区分：严格归约、可证明上界、启发式指标和仅用于建立直觉的类比。

\section{Sketch 是受限长度的消息}

数据流算法把长输入流压缩成小状态 $S$。若状态使用 $s$ bits，那么它最多有 $2^s$ 种取值。查询算法只能依据 $S$ 和查询参数输出答案。

这与单向通信协议非常接近：
\begin{enumerate}
  \item Alice 读取流的前一部分并生成 sketch $S$；
  \item Alice 把 $S$ 作为消息发送给 Bob；
  \item Bob 根据自己的输入继续更新或发起查询；
  \item 若 sketch 能解决流问题，通信协议就能解决相应通信问题。
\end{enumerate}
因此，如果相关通信问题需要至少 $b$ bits，任何实现该归约的 sketch 也必须使用至少 $b$ bits。

\begin{intuitionbox}
流式空间下界常问：有限状态究竟区分了多少种输入情况？信息论提供“状态容量”的语言，通信复杂度则提供把空间限制转成消息长度限制的标准模型。
\end{intuitionbox}

\section{从熵界到通信下界的基本骨架}

设 Alice 的随机输入为 $X$，消息为 $M=f(X)$，长度至多 $s$ bits，则
\[
H(M)\le s,
\qquad
I(X;M)\le H(M)\le s.
\]
若任务要求 Bob 从 $(M,Y)$ 中以小错误恢复 $X$ 的某个高信息量函数，就可借助 Fano 不等式或问题特定论证，推出 $I(X;M|Y)$ 必须足够大，进而得到 $s$ 的下界。

真正的下界证明还需要：
\begin{itemize}
  \item 选择困难输入分布；
  \item 明确允许的随机性（public/private coins）；
  \item 明确单向还是多轮通信；
  \item 把流算法的成功概率、passes 和空间逐项映射到协议。
\end{itemize}

\section{预测作为 side information}

在 learning-augmented streaming 中，预测 $Z$ 可以被看作关于真实流特征 $X$ 的 side information。自然问题包括：
\begin{align*}
&H(X|Z)\text{ 比 }H(X)\text{ 小多少？}\\
&I(X;Z)=H(X)-H(X|Z)\text{ 是否能转化为空间收益？}
\end{align*}

Slepian--Wolf 理论表明，在特定的独立同分布分布式无失真压缩模型中，解码端拥有 $Z$ 时，$X$ 的渐近压缩率可下降到 $H(X|Z)$。但学习增强 Sketch 往往是 worst-case 流、有限内存、在线更新和查询误差模型，不能直接套用该结论。

\begin{warningbox}
预测的 KL 误差小，不自动推出某个 sketch 的查询误差小；互信息高，也不自动给出高效在线算法。必须提供把信息量连接到具体状态、更新和查询保证的定理。
\end{warningbox}

\section{Consistency 与 robustness}

学习增强算法通常同时要求：
\begin{description}
  \item[Consistency] 预测准确时优于经典基线；
  \item[Robustness] 预测任意错误时不会比经典基线坏太多。
\end{description}
一个完整问题必须定义预测接口 $Z$、误差度量 $\eta(X,Z)$、资源成本和查询误差。例如：
\[
\text{space}=S(\epsilon,\delta,\eta),
\qquad
\Prb[|\widehat f-f|>E(\eta)]\le\delta.
\]
信息论可以帮助设计 $\eta$ 或证明不可能性，但最终保证要落到算法参数上。

\section{Verifier 反馈包含多少信息}

在 verifier-guided search 中，候选对象为 $C$，反馈为 $F$。若 verifier 只返回一位 accept/reject，则单次反馈熵至多 1 bit；若返回详细反例、proof state 或连续分数，理论上可能携带更多关于候选错误结构的信息。

但“反馈字段更多”不等于“搜索有效信息更多”：
\begin{itemize}
  \item 反馈可能高度冗余；
  \item 搜索器可能无法利用其结构；
  \item 反馈可能泄漏 benchmark 答案并导致投机；
  \item 生成反馈的计算成本可能远高于收益。
\end{itemize}
可以把研究问题写成：在固定 verifier 成本下，哪种反馈表示提高了候选成功率或降低了搜索样本复杂度？

\section{证明轨迹、程序与描述长度}

证明、程序和算法都可以编码成有限字符串。描述长度可以用来提出问题：
\begin{itemize}
  \item 一个候选是否只是现有解的冗长等价变换？
  \item 搜索空间的表示是否引入大量无意义自由度？
  \item 证明证书能否被短消息验证？
\end{itemize}
但不能把“描述更短”直接等同于“数学更深”或“算法复杂度更优”。Kolmogorov complexity、最小描述长度和 Shannon 熵属于不同层次的对象，需要各自的模型和假设。

\section{适合当前背景的研究问题}

\begin{aibox}
一个可控的两周问题可以这样定义：
\begin{enumerate}
  \item 选择 Count-Min Sketch 作为经典 baseline；
  \item 令预测器给出候选 heavy hitters 或频率先验；
  \item 明确定义预测误差，而不是只给模型准确率；
  \item 比较 exact counter、经典 sketch 和预测增强版本；
  \item 记录空间、查询误差、失败概率和最坏预测分布；
  \item 只把实验证据写成经验结果，把 consistency/robustness 留给可证明部分。
\end{enumerate}
\end{aibox}

\section{证据等级表}

\begin{center}
\begin{tabularx}{\textwidth}{lXX}
\toprule
证据 & 可以说明 & 不能单独说明\\
\midrule
有限数据实验 & 特定分布和预算上的表现 & worst-case 或渐近保证\\
随机测试 & 抽样范围内未发现错误 & 任意输入正确\\
通信归约 & 模型内的空间/通信下界 & 其他模型仍然不可能\\
信息不等式 & 随机变量间的必然约束 & 存在高效实现算法\\
Lean/kernel 检查 & 当前形式化命题可证明 & 自然语言语义正确或有创新性\\
\bottomrule
\end{tabularx}
\end{center}

\section{小结与验收}

\begin{checkpointbox}
\begin{enumerate}
  \item 画出一次从流算法到单向通信协议的归约流程；
  \item 解释消息长度 $s$ 为什么给出 $I(X;M)\le s$；
  \item 给出一个 side information 有用但不能直接套 Slepian--Wolf 的流式例子；
  \item 为一个 learning-augmented sketch 写出 prediction interface、error、consistency 和 robustness；
  \item 区分 verifier 反馈的字段长度、熵和对搜索真正有用的信息。
\end{enumerate}
\end{checkpointbox}

\chapter{综合习题与参考答案}

建议先独立完成，再查看参考答案。答案强调关键步骤，不替代完整书写。

\section{习题}

\begin{enumerate}[label=\textbf{I\arabic*.}]
  \item 计算分布 $(1/2,1/4,1/8,1/8)$ 的熵，并说明为什么它能被码长 $(1,2,3,3)$ 精确达到。
  \item 设 $X\sim\mathrm{Bernoulli}(p)$。证明 $h_2(p)$ 关于 $p=1/2$ 对称且严格凹，并找出最大值。
  \item 给出一个非奇异但不可唯一译的码，再给出一个唯一可译但非前缀码的码（可查阅 Sardinas--Patterson 判据验证后者）。
  \item 判断码长集合 $(1,2,3,3)$、$(1,2,2,3)$、$(2,2,3,3,3)$ 是否满足 Kraft 不等式。
  \item 对分布 $(0.4,0.3,0.2,0.1)$ 构造 Huffman 码，计算平均码长并与熵比较。
  \item 设 $Y=X\oplus N$，其中 $X\sim\mathrm{Bernoulli}(1/2)$、$N\sim\mathrm{Bernoulli}(q)$ 且独立。计算 $H(Y)$、$H(Y|X)$ 和 $I(X;Y)$。
  \item 证明若 $X\to Y\to Z$ 是 Markov 链，则 $I(X;Z)\le I(X;Y)$。
  \item 解释固定长度信源编码中，为什么 $R>H(X)$ 时典型集可以被索引，而 $R<H(X)$ 时不够。
  \item 推导 BEC$(\varepsilon)$ 的容量，并比较它和 BSC$(\varepsilon)$ 的数值及语义差别。
  \item 在随机编码证明中，若 $R=I(X;Y)-2\delta$，估计错误码字联合典型事件的 union bound 指数阶。
  \item 使用 Fano 不等式解释：若从 $s$-bit 消息中恢复一个均匀 $m$-bit 字符串且错误概率趋零，为什么必须有 $s\gtrsim m$。
  \item 为“预测增强 Count-Min Sketch”写一份信息论问题定义，至少包含随机变量、预测误差、资源、成功事件和不声称的结论。
\end{enumerate}

\section{参考答案与提示}

\paragraph{I1.}
自信息为 $1,2,3,3$，所以 $H=1.75$ bits。概率满足 $p_i=2^{-\ell_i}$，且 Kraft 和等于 1，因此整数码长与理想自信息完全一致。

\paragraph{I2.}
直接代换 $1-p$ 得对称性。二阶导
\[
h_2''(p)=-\frac1{\ln2}\left(\frac1p+\frac1{1-p}\right)<0
\]
给出严格凹；一阶导在 $p=1/2$ 为零，最大值为 1。

\paragraph{I3.}
非奇异但不可唯一译可用 $\{0,01,10\}$，因为 $010$ 有两种解析。唯一可译但非前缀码可用 $\{0,01\}$：$0$ 是 $01$ 的前缀，但任何合法串中出现的 1 只能与它前面的 0 组成 $01$，因此解析唯一。严格处理更大码集时可使用 Sardinas--Patterson 判据，不能只凭“短时间没找到歧义”下结论。

\paragraph{I4.}
分别计算 Kraft 和：
\[
\tfrac12+\tfrac14+\tfrac18+\tfrac18=1,
\quad
\tfrac12+\tfrac14+\tfrac14+\tfrac18=\tfrac98>1,
\quad
\tfrac14+\tfrac14+3\cdot\tfrac18=\tfrac78.
\]
第一、第三组可构造前缀码，第二组不行。

\paragraph{I5.}
合并 $0.1+0.2=0.3$，再合并两个 $0.3$ 得 $0.6$，最后 $0.4+0.6$。一种码长为 $(1,2,3,3)$，平均码长 $1.9$ bits。熵约为 $1.846$ bits，满足 $H\le L<H+1$。

\paragraph{I6.}
$Y$ 均匀，所以 $H(Y)=1$；给定 $X$ 后只剩噪声，$H(Y|X)=h_2(q)$；因此 $I(X;Y)=1-h_2(q)$。

\paragraph{I7.}
用互信息链式法则：
\[
I(X;Y,Z)=I(X;Y)+I(X;Z|Y)=I(X;Y),
\]
同时
\[
I(X;Y,Z)=I(X;Z)+I(X;Y|Z)\ge I(X;Z).
\]

\paragraph{I8.}
典型集约有 $2^{nH}$ 个元素。码率 $R$ 提供 $2^{nR}$ 个索引。$R>H$ 时足够覆盖高概率典型集；$R<H$ 时可编码的典型序列比例按 $2^{-n(H-R)}$ 衰减。

\paragraph{I9.}
擦除位置以概率 $\varepsilon$ 不提供信息，容量为 $1-\varepsilon$。BSC 容量为 $1-h_2(\varepsilon)$；擦除位置可见，翻转位置不可见，因此两者不能只按“受损比例”比较。

\paragraph{I10.}
错误码字数约 $2^{nR}$，单个伪匹配概率约 $2^{-n(I-\delta)}$，乘积为
\[
2^{-n(I-R-\delta)}=2^{-n\delta}.
\]

\paragraph{I11.}
设原串为 $X$、消息为 $M$、恢复结果为 $\hat X$。Fano 给出 $H(X|M)=o(m)$，于是
\[
m=H(X)=I(X;M)+H(X|M)\le H(M)+o(m)\le s+o(m).
\]

\paragraph{I12.}
一种合格定义是：$X$ 为真实频率向量，$Z$ 为预测向量，误差 $\eta=\|X-Z\|_1/\|X\|_1$；状态限于 $s$ bits；成功事件为所有查询误差满足给定界，失败概率至多 $\delta$。实验只能说明选定分布上的经验表现，不能单独声称 worst-case consistency/robustness。


\begin{thebibliography}{9}
\bibitem{cover-thomas}
Thomas M. Cover and Joy A. Thomas.
\newblock \emph{Elements of Information Theory}, 2nd edition.
\newblock Wiley, 2006.

\bibitem{mackay}
David J. C. MacKay.
\newblock \emph{Information Theory, Inference, and Learning Algorithms}.
\newblock Cambridge University Press, 2003.

\bibitem{yeung}
Raymond W. Yeung.
\newblock \emph{A First Course in Information Theory}.
\newblock Springer, 2002.

\bibitem{cormode-yi}
Graham Cormode and Ke Yi.
\newblock \emph{Small Summaries for Big Data}.
\newblock Cambridge University Press, 2020.
\end{thebibliography}

\end{document}
